\begin{tikzpicture}[
  font=\sffamily,
  stage/.style={
    rectangle, rounded corners=3.5pt, draw=ceindigo, line width=0.95pt,
    fill=cepanel, text width=2.85cm, align=center, minimum height=2.15cm,
    inner sep=5pt
  },
  support/.style={
    rectangle, rounded corners=3pt, draw=ceteal, line width=0.8pt,
    fill=cepanelteal, text width=4.0cm, align=center, minimum height=1.35cm,
    inner sep=6pt
  },
  risk/.style={
    rectangle, rounded corners=2pt, draw=cecrimson!80, line width=0.7pt,
    fill=cepanelcrimson, text width=3.05cm, align=center, minimum height=1.12cm,
    inner sep=5pt
  },
  arrow/.style={->, >=Stealth, line width=1.2pt, ceindigo}
]
  \node[font=\sffamily\bfseries\small, ceindigo] (t1) at (5.25, 2.35)
    {LAS CUATRO FASES DEL PROCESO EMPRENDEDOR (HISRICH, PETERS Y SHEPHERD)};

  \node[stage] (f1) at (0, 0.45) {
    \textcolor{cecrimson}{\bfseries\scriptsize FASE 1}\\[2pt]
    \textbf{\footnotesize Identificación}\\[3pt]
    \scriptsize\color{ceslate} Alerta emprendedora, evaluación de dolor real y ventana temporal.
  };
  \node[stage, right=0.50cm of f1] (f2) {
    \textcolor{cecrimson}{\bfseries\scriptsize FASE 2}\\[2pt]
    \textbf{\footnotesize Plan de empresa}\\[3pt]
    \scriptsize\color{ceslate} Diseño de modelo, validación y plan formal de viabilidad.
  };
  \node[stage, right=0.50cm of f2] (f3) {
    \textcolor{cecrimson}{\bfseries\scriptsize FASE 3}\\[2pt]
    \textbf{\footnotesize Recursos}\\[3pt]
    \scriptsize\color{ceslate} Inventario propio, captación de capital, alianzas y alta legal.
  };
  \node[stage, right=0.50cm of f3] (f4) {
    \textcolor{cecrimson}{\bfseries\scriptsize FASE 4}\\[2pt]
    \textbf{\footnotesize Gestión}\\[3pt]
    \scriptsize\color{ceslate} Arranque operativo, sistemas de control y escalado sostenible.
  };

  \draw[arrow] (f1.east) -- (f2.west);
  \draw[arrow] (f2.east) -- (f3.west);
  \draw[arrow] (f3.east) -- (f4.west);

  % Bloque intermedio: los 4 riesgos de Eisenmann
  \node[font=\sffamily\bfseries\footnotesize, cecrimson] at (5.25, -1.25)
    {LOS CUATRO RIESGOS SIMULTÁNEOS A GESTIONAR (EISENMANN)};

  \node[risk] (r1) at (0.05, -2.18)  {\textbf{\scriptsize Riesgo de demanda}\\[-1pt]\scriptsize ¿Comprará el cliente?};
  \node[risk] (r2) at (3.52, -2.18)  {\textbf{\scriptsize Riesgo tecnológico}\\[-1pt]\scriptsize ¿Es técnicamente viable?};
  \node[risk] (r3) at (6.98, -2.18)  {\textbf{\scriptsize Riesgo financiero}\\[-1pt]\scriptsize ¿Habrá caja suficiente?};
  \node[risk] (r4) at (10.45, -2.18) {\textbf{\scriptsize Riesgo de ejecución}\\[-1pt]\scriptsize ¿Cumplirá el equipo?};

  % Bloque inferior: las 3 redes de apoyo
  \node[font=\sffamily\bfseries\footnotesize, ceteal!90!black] at (5.25, -3.35)
    {ECOSISTEMA DE SOPORTE: LAS TRES REDES DE APOYO DEL EMPRENDEDOR};

  \node[support] (s1) at (0.65, -4.35) {
    \textbf{\footnotesize Red Moral (Personal)}\\[2pt]
    \scriptsize\color{ceslate} Familia y entorno cercano. Apoyo emocional ante la incertidumbre.
  };
  \node[support] (s2) at (5.25, -4.35) {
    \textbf{\footnotesize Red Profesional}\\[2pt]
    \scriptsize\color{ceslate} Mentores, socios, asesores, primeros clientes y asociaciones.
  };
  \node[support] (s3) at (9.85, -4.35) {
    \textbf{\footnotesize Red Institucional}\\[2pt]
    \scriptsize\color{ceslate} UGR, CADE / Andalucía Emprende, ENISA, viveros y PAE.
  };
\end{tikzpicture}